\documentclass[11pt,letterpaper]{article}
\usepackage[T1]{fontenc}
\usepackage[utf8]{inputenc}
\usepackage{lmodern,microtype}
\usepackage[margin=1in]{geometry}
\usepackage{amsmath,amssymb,amsthm,mathtools}
\usepackage{booktabs,longtable,array}
\usepackage[dvipsnames]{xcolor}
\usepackage{enumitem,fancyhdr}
\usepackage{listings}
\usepackage[colorlinks=true,linkcolor=MidnightBlue,urlcolor=MidnightBlue,citecolor=MidnightBlue,breaklinks=true]{hyperref}
\usepackage{xurl}
\hypersetup{pdftitle={Smaller direct quartic equations for ZFC consistency},pdfauthor={Research assistance for Vladimir Reshetnikov}}
\newtheorem{theorem}{Theorem}[section]
\newtheorem{lemma}[theorem]{Lemma}
\newtheorem{corollary}[theorem]{Corollary}
\theoremstyle{definition}
\newtheorem{definition}[theorem]{Definition}
\newtheorem{remark}[theorem]{Remark}
\newcommand{\Z}{\mathbb Z}
\newcommand{\N}{\mathbb N}
\newcommand{\Ftwo}{\mathbb F_2}
\newcommand{\Con}{\operatorname{Con}}
\newcommand{\ZFC}{\mathsf{ZFC}}
\newcommand{\ZF}{\mathsf{ZF}}
\newcommand{\supp}{\operatorname{supp}}
\newcommand{\sgn}{\operatorname{sgn}}
\newcommand{\degp}{\operatorname{deg}}
\newcommand{\MainVars}{3776}
\newcommand{\MainLast}{3775}
\newcommand{\MainTerms}{20194}
\newcommand{\MainRows}{2015}
\newcommand{\MainLastRow}{2014}
\newcommand{\MainHeight}{47959}
\newcommand{\MainSaved}{4843}
\newcommand{\MainPercent}{56.19}
\newcommand{\ElimVars}{3282}
\newcommand{\AffineElims}{1915}
\newcommand{\QuadraticElims}{1367}
\newcommand{\MainPolynomialHash}{28e218143296c666cd387ef4d9ccc5a4ab0d45506a59ef2fb946c37291a17337}
\newcommand{\MainVarsComma}{3,776}
\newcommand{\MainRowsComma}{2,015}
\newcommand{\MainTermsComma}{20,194}
\newcommand{\MainHeightComma}{47,959}

\setlength{\emergencystretch}{3em}
\setlist{itemsep=3pt,topsep=5pt}
\lstset{basicstyle=\ttfamily\footnotesize,breaklines=true,columns=fullflexible,keepspaces=true,frame=single,rulecolor=\color{black!20}}
\pagestyle{fancy}
\fancyhf{}
\fancyhead[L]{\small Direct quartic consistency equations}
\fancyhead[R]{\small ProveIt research supplement}
\fancyfoot[C]{\thepage}
\setlength{\headheight}{14pt}
\title{Smaller direct quartic equations for ZFC consistency\\[0.3em]\large Signed witnesses, sign normalization, and polynomial elimination}
\author{Research supplement prepared for Vladimir Reshetnikov}
\date{23 October 2026\\\small Repository snapshot \texttt{278bc68df1af}}

\begin{document}
\maketitle
\begin{abstract}
We construct an explicit integer-coefficient polynomial $P$ of exact degree
four in $\MainVars$ unrestricted integer variables. On the established
ordinary mathematical dependencies of ProveIt's direct C222 construction,
ZFC proves that $P$ has an integer zero exactly when a finite syntactic ZFC
refutation exists. The complete polynomial has $\MainTerms$ nonzero
monomials and coefficient height $\MainHeight$. Its zero set is specified
both by a fully expanded coefficient file and by $\MainRows$ quadratic
residuals whose squares sum to $P$.

The starting direct quartic used 8,619 integers. The new reduction removes
an unnecessary natural-number encoding of 221 signed quotients, normalizes
110 signs using exact polynomial symmetries, derives hundreds of other
sign conditions from existing equations, and eliminates auxiliary
coordinates with polynomial inverse maps. The resulting dimension
reduction is $\MainSaved$ variables, or $\MainPercent\%$, with no increase
in degree, while the coefficient height drops from 127,285,349 to
47,959. A second equation uses only 3,671 variables at the larger height
3,835,369,234. All new transformations have universal
mathematical proofs and finite exact certificates checked by independent
programs. This is a measured improvement of the direct, moderate-coefficient
quartic family, not a minimum-variable record or a proof of global
simplicity. No new end-to-end Lean proof or explicitly emitted
object-language ZFC derivation is claimed.
\end{abstract}

\tableofcontents
\newpage

\section{The result and the meaning of simplicity}
\label{sec:result}

The task is to give a fixed polynomial whose lack of integer solutions
expresses syntactic consistency of ZFC, with a construction simpler than
the existing alternatives. There is no agreed total ordering of
polynomials by simplicity. Number of variables, degree, number of terms,
coefficient height, description length, and complexity of a correctness
proof can favor different constructions. A claim of the globally
``simplest'' equation would need both a specified ordering and an
optimality argument. We make neither claim.

This supplement improves a particular useful region of that comparison:
\emph{explicit quartics with ordinary literal coefficients, obtained
directly from a finite proof calculus}. The starting equation is the
8,619-integer companion of the 3,609-natural-variable shared-CRT prime-lift
construction at repository commit
\texttt{278bc68df1af125d62f9cbfe28c9fd3ebcf50d5b}~\cite{repo-integer,repo-prime}.
Its arithmetic represents C222 derivations and uses neither the K215
Turing machine nor Jones's universal Diophantine expression.

\begin{theorem}[The explicit new equation]
\label{thm:main}
Let $P\in\Z[X_0,\ldots,X_{\MainLast}]$ be the polynomial in the delivered
\texttt{equations/main/expanded\_polynomial.json}, equivalently the sum of
the squared residuals in the adjacent \texttt{equation.json}. Then
\[
\boxed{\quad
\ZFC\vdash
\left[\neg\exists\boldsymbol x\in\Z^{\MainVars}\;
P(\boldsymbol x)=0\right]\ \longleftrightarrow\ \Con(\ZFC).
\quad}
\]
Here $\Con(\ZFC)$ means that no finite derivation of falsity exists from
the repository's specified sealed ZFC axiom predicate. The ordinary proof
uses the named C222/calculus and finite arithmetic representation theorems
listed in Section~\ref{sec:source}, together with the new transformations
proved below. The polynomial has the exact measurements
\[
\begin{aligned}
\text{unknowns}&=\MainVars,&
\text{degree}&=4,\\
\text{nonzero monomials}&=\MainTerms,&
\max|\text{coefficient}|&=\MainHeight.
\end{aligned}
\]
Every unknown ranges over \emph{all} integers. There are no free
parameters or unspecified coefficient choices.
\end{theorem}

The notation $\ZFC\vdash$ states ordinary mathematical provability. The
proof consists of uniform finite constructions, elementary arithmetic,
classical number theory, and the named ordinary logical calibration
theorems. It does not assert that this archive contains a new
machine-checked derivation of that sentence. In particular, the proof
does not infer ZFC-provability merely from the existence of a Lean
declaration. The distinction is discussed in Section~\ref{sec:internal}.

\subsection{Measured improvements}

\begin{table}[ht]
\centering\small
\begin{tabular}{lrrrr}
\toprule
Construction & Integer & Quadratic & Quartic & Coefficient\\
& unknowns & residuals & terms & height\\
\midrule
Frozen direct parent & 8,619 & 4,619 & 39,811 & 127,285,349\\
Signed quotients & 6,628 & 3,733 & 31,237 & 127,285,349\\
Reduced sign guards & 4,879 & 3,150 & 22,910 & 127,285,349\\
Retained-wire guards & 7,058 & 5,297 & 24,277 & 1,992\\
Fewer-variable companion & 3,671 & 1,942 & 26,271 & 3,835,369,234\\
Recommended $P$ & \MainVarsComma & \MainRowsComma & \MainTermsComma & \MainHeightComma\\
\bottomrule
\end{tabular}
\caption{Complete materialized polynomials, all over unrestricted
integers and all of exact degree four. ``Height'' refers to the expanded
polynomial, not just the residual coefficients.}
\label{tab:main}
\end{table}

Two source presentations give different useful final equations. More
aggressive elimination of the affine-source intermediate reaches 3,671
variables, but increases coefficient height substantially. Starting
from the retained-wire presentation and limiting substitutions to
two-term right-hand sides instead yields the recommended
3,776-variable polynomial. It has only 105 more variables, 6,077 fewer
monomials, and approximately 79,972 times smaller coefficient height.
It improves all tabulated size measures of the frozen 8,619-variable
parent, with degree fixed at four. The 7,058-variable equation remains
available when minimizing coefficient height itself is the priority.

Small-variable alternatives already exist. The repository has nine
natural-variable and ten integer-variable equations of much higher
degree; its latest accepted ten-integer table includes degree 1,080,888
at 512 arithmetic operations~\cite{repo-frontier}. Sun's 2026 preprint
gives a general ten-integer-variable representation theorem~\cite{sun2026}.
The present work does not improve that variable count. Likewise,
proof-theoretic Diophantine encodings independent of a Turing-machine
simulation are not a new general idea; Carl and Moroz developed such an
approach explicitly~\cite{carl-moroz}. Our contribution is the concrete
reduction of a fully specified direct quartic and the reusable exact
transformation certificates.

\section{A literal polynomial, not an arithmetic shorthand}
\label{sec:literal}

All construction variables are enumerated in a fixed list. In a sparse
row, an entry $(c,[i_1,\ldots,i_s])$ denotes the ordinary monomial
$cX_{i_1}\cdots X_{i_s}$; the empty list denotes the constant $c$.
Repeated indices denote ordinary integer powers. Every residual in the
final equation has $s\le2$.

Write the final residual list as $f_0,\ldots,f_{\MainLastRow}$. The defining
equation is simply
\begin{equation}
\label{eq:definition}
\boxed{\qquad P(\boldsymbol X)
=\sum_{r=0}^{\MainLastRow} f_r(\boldsymbol X)^2=0,
\qquad\boldsymbol X\in\Z^{\MainVars}.\qquad}
\end{equation}
Both every $f_r$ and every collected coefficient of $P$ are provided.
There are no variable exponents, factorials, binomial coefficients,
bounded universal quantifiers, remainder operators, or residual
computability predicates in these files. Such operations occur only in
the proof explaining how the earlier natural parent was constructed.

Over the integers, a sum of squares vanishes exactly when each summand
vanishes. Thus the residual presentation and the expanded polynomial
have identical zero sets. This is the only conjunction operation in the
final definition. The upper bound $\deg P\le4$ follows immediately;
the exact expansion contains nonzero degree-four terms, proving equality.
The independent checker also verifies that every declared variable
occurs in the expanded polynomial.

The JSON format is not an additional mathematical parameter. It is an
unambiguous finite coefficient table, in the same sense as a printed
list of integers. The source, construction algorithms, certificates,
and hashes make the table reproducible. A human-readable residual
listing is included for inspection without a JSON viewer.

\section{The established direct source and its logical meaning}
\label{sec:source}

This section identifies the imported results and summarizes the
unbounded source argument. It is included to prevent a local polynomial
optimization from being mistaken for an independent construction of the
entire consistency encoding.

\subsection{The exact natural parent}

Let $\mathcal S_0$ be the frozen natural system in
\texttt{signed\_quotients/source/natural\_equation.json}. Its complete byte string has
SHA-256
\begin{center}\small\ttfamily
e851ff1d420a39be069e63f6f117557356c38e6d8272a481258fe38cb38dc6ea.
\end{center}
It has 3,609 nonnegative coordinates and 2,949 residuals of degree at most
two. Its squared sum has 15,752 monomials and height 127,285,349. Of its
coordinates, 1,670 are semantic inputs and 1,939 are retained arithmetic
wires. The parent generator is bundled and reconstructs the complete
equation byte for byte. Its arithmetic calls expand 55 positive power
graphs, 16 binomial graphs, and two factorial graphs into ordinary
quadratic rows.

The same source also has a retained-wire natural presentation with
6,180 variables, 5,520 quadratic rows, 17,421 quartic monomials, and
coefficient height 1,992. Eliminating 2,571 nonnegative affine wire
definitions yields exactly $\mathcal S_0$. Each deleted wire has a
unique nonnegative polynomial restoration, so the two natural
presentations have equivalent zero existence. Both complete parents
are regenerated by the supplied source verifier. The retained
presentation, whose literal file hash begins \texttt{39b3e264c051},
is the starting point of the recommended low-height final equation.

The established ordinary source theorem is
\begin{equation}
\label{eq:parent-correctness}
\exists\boldsymbol y\in\N^{3609}\ \mathcal S_0(\boldsymbol y)
\quad\longleftrightarrow\quad
\text{there exists a finite syntactic ZFC refutation}.
\end{equation}
We import the proved source and representation theorems behind this
statement, rather than assume a new simulator calibration. The current
repository distinguishes this direct chain from some native
machine/coefficient chains whose complete formal bridges remain separate
research tasks~\cite{repo-direct-review,repo-shared-review}.

\subsection{Numerical calculus}

The numerical calculus C222 has truth label $0$, falsity label $1$, a
fixed finite set of numerical rule families, modus ponens, and
generalization. Its constructors are ordinary polynomials derived from
Cantor pairing; for example,
\[
\operatorname{Imp}(a,b)=2(a+b)(a+b+1)+4b.
\]
The relation, quantifier, and equality constructors use fixed offsets.
The supplied direct source includes all twelve padded single-conclusion
branches. Unused premise positions receive the truth label $0$.

The decisive source declarations are pinned to commit
\texttt{9a33342a03d1} (the full revision is in the source ledger):
\begin{lstlisting}
KEndpoint.c222_iff_zf_refutation :
  KCode222.C222 Code254.F <-> BProv ZFax_s [] Form.fBot

ConstructibleChoice.zf_refutation_iff_zfc_refutation :
  BProv ZFax_s [] Form.fBot <-> BProv ZFCax_s [] Form.fBot
\end{lstlisting}
These are unconditional calculus/refutation interfaces, not
hypotheses about execution of a Turing machine. The ordinary arguments
use first-order soundness and completeness, finite proof transformations,
and the constructible-model relative-consistency theorem. They do not
require an externally transitive model or an assumption of
$\Con(\ZFC)$. We retain the source pins and attribution and do not claim
to have rebuilt the complete Lean dependency closure.

\subsection{Finite flow and the product carrier}

For finitely many valid rule instances let $c_i$ be the conclusion and
$p_i,q_i$ their truth-padded premises. The balance at every positive
label $j$ is
\begin{equation}
\label{eq:flow}
\#\{i:c_i=j\}=[j=1]+\#\{i:p_i=j\}+\#\{i:q_i=j\}.
\end{equation}
A proof tree gives these balances. Conversely, take the closure of truth
under the finitely supported rules. If $1$ were outside this closure,
sum the conclusion-minus-premise balance over labels outside it. Every
rule whose conclusion is outside has at least one premise outside, so
its contribution is nonpositive. Rules whose conclusion is inside also
contribute at most zero. The target contributes $1$, a contradiction.
This argument allows cycles, repeated premises, and repeated rule
instances. It uses one conclusion per rule and nonnegative
multiplicities.

If there are $t\ge1$ instances with labels at most $m\ge1$, define
\[
A(Z)=\prod_{i<t}(1+c_iZ),\qquad
B(Z)=(1+Z)\prod_{i<t}(1+p_iZ)(1+q_iZ),\qquad M=(m+1)^t.
\]
Equality $A(M)=B(M)$ first bounds the number of nonconstant factors of
$B$ by $t$: otherwise $B(M)>M^{t+1}$, whereas
$A(M)<M^{t+1}$. Consequently $A(1),B(1)\le M$, and each nonconstant
coefficient is strictly less than $M$, because the constant coefficient
is one. A least nonzero coefficient of $A-B$, reduced modulo $M$, would
be both divisible by $M$ and smaller than $M$ in absolute value. Thus
$A=B$ as polynomials. Multiplicities of the distinct roots $-1/j$ give
\eqref{eq:flow}.

A single accumulator encodes the equality by positive transitions
\[
S_0=M+1,\qquad S_t=1,\qquad
S_{i+1}(1+c_iM)=S_i(1+p_iM)(1+q_iM).
\]
A proof tree supplies integral accumulator values by following its open
obligations. The accumulator sequence is beta-coded with moduli
$1+(i+1)y$. The proof, label bound, carrier, trace, beta code, and later
row-witness bound are chosen in that order. The label bound and the
bound on all arithmetic witnesses are distinct, so there is no circular
size requirement.

\subsection{Uniform CRT and prime-power lifting}

The bounded row relation has nonnegative-coefficient polynomial sides
$E,F$, an index bound $N$, external parameters $z$, and coordinate
bounds $u_1,\ldots,u_d$. Put
\[
H=E(N,z,u)+F(N,z,u)+N,\quad T=H!,\quad
D_i=1+(i+1)T,\quad B=\prod_{i<N}D_i.
\]
The natural source uses encoded columns $A_j$ and encoded index $C$
satisfying
\begin{align}
B&\mid\binom{A_j}{u_j+1} &&(1\le j\le d),\label{eq:crt-range}\\
B&\mid1+(C+1)T,\qquad
E(C,z,A)\equiv F(C,z,A)\pmod B.\label{eq:crt-row}
\end{align}
These conditions are equivalent to the existence of a bounded valid row
for every $i<N$~\cite{repo-prime}.

For soundness choose a prime $p\mid D_i$. It is coprime to $H!$, hence
$p>H$. Multiplying \eqref{eq:crt-range} by $(u_j+1)!$ and using the
falling-factorial identity shows that $p$ divides some $A_j-r$ with
$0\le r\le u_j$. This remains true when $A_j\le u_j$, since then the
falling product has a zero factor. Therefore each column has a bounded
decoded representative modulo $p$. The index congruence yields $C=i$
modulo $p$. Both decoded values $E,F$ are at most $H<p$, so their
congruence is equality.

For completeness the $D_i$ are pairwise coprime. For each coordinate
let $R_j=(u_j+1)!$ and let $R_{ij}$ contain precisely the prime-power
part of $R_j$ supported on primes dividing $D_i$. Choose the column by
CRT with
\[
A_j\equiv w_j^{(i)}\pmod{D_iR_{ij}},\qquad
C\equiv i\pmod{D_i}.
\]
These lifted column moduli remain pairwise coprime.
The falling product is divisible by $D_iR_{ij}$; cancel $R_{ij}$ and
then the coprime remaining factor $R_j/R_{ij}$ to obtain
\eqref{eq:crt-range}. A positive column is chosen modulo the product of
the lifted periods. Merely adding $B$ is not sufficient without
factorial coprimality: for instance, $3\mid\binom13$ but
$3\nmid\binom43$.

This construction does not require $H\ge u_j+1$. All factorization and
CRT operations here construct witnesses in the proof; they are not
operators left in the final polynomial. The supplied finite Pell-based
power, binomial, and factorial representations replace the displayed
arithmetic macros by ordinary quadratic rows. Their exact call-site
positivity hypotheses and coefficient identities are audited by the
bundled original independent checkers.

The central arithmetic interfaces are
\texttt{PellRepresentation.criterion\_iff\_pell} and
\texttt{PowerRepresentation.criterion\_iff\_pow}. Their source uses
the finite Matiyasevich Pell criterion and its power consequence,
respectively; selected exact source files are retained under
\texttt{source\_calibration}. Appendix~\ref{app:power} displays the
particular positive power criterion used here. This is a named
number-theoretic dependency, not an unexpanded power operation in
the final coefficient table.

\section{Removing the natural encoding of signed quotients}
\label{sec:quotients}

The first new reduction changes the domain of a witness to the domain
the mathematical relation actually needs. An integer congruence
quotient may be negative. The natural parent represents it by a
difference of two nonnegative witnesses, then computes the two products
separately. Applying a generic natural-to-integer conversion to those
witnesses adds more unnecessary variables.

\begin{lemma}[Private signed quotient]
\label{lem:signed}
Suppose a natural system contains private coordinates $q_+,q_-,u,v$
used only in
\[
u=Mq_+,\qquad v=Mq_-,\qquad E+c(u-v)=0,
\]
where $M$ is a nonnegative expression in retained natural coordinates,
$c\in\Z$, and $E$ is independent of these four coordinates. Replace the
three rows by
\[
E+cMz=0,\qquad z\in\Z.
\]
Existence of a solution of the original natural system is equivalent to
existence of a solution of the new mixed-domain system.
\end{lemma}

\begin{proof}
Projection takes $z=q_+-q_-$. Conversely, from a signed witness $z$
choose
\[
q_+=\max(z,0),\quad q_-=\max(-z,0),\quad
u=Mq_+,\quad v=Mq_-.
\]
These are natural numbers and satisfy all three original rows. Private
use ensures that no other constraint has changed. The proof also covers
$M=0$ and does not divide by the modulus.
\end{proof}

The replacement is not a bijection: many pairs $(q_+,q_-)$ have the
same difference. This is immaterial for existential equivalence, and
the displayed canonical lift is sufficient. It would be incorrect to
claim that simply allowing both old quotient coordinates to be signed
and discarding their guards leaves an arbitrary natural source sound;
the private-use property and the explicit lift are essential.

In the frozen parent there are exactly 221 qualifying pairs. The
certificate checks the complete occurrence set of each pair, its common
affine modulus, its two private product wires, and their opposite linear
coefficients in the single final use row. Two additional private
affine bridge wires are first substituted out. Each quotient block
then removes three coordinates and two defining rows. Thus
\[
3609-3\cdot221-2=2944,\qquad
2949-2\cdot221-2=2505.
\]
The resulting core has 221 signed quotient coordinates, 1,228 natural
semantic coordinates, and 1,495 natural arithmetic wires.

All retained gate definitions have nonnegative coefficients in retained
natural coordinates and contain no signed quotient. The common moduli
also have this property. Therefore the restoration in
Lemma~\ref{lem:signed} genuinely returns natural product wires; this is
verified syntactically rather than assumed from variable names.

\subsection{The initial integer lift}

For completeness, the basic nonnegativity certificate is
\begin{equation}
\label{eq:three-square}
x=a^2+b^2+c(c+1),\qquad a,b,c\in\Z.
\end{equation}

\begin{lemma}
An integer $x$ admits \eqref{eq:three-square} if and only if $x\ge0$.
\end{lemma}
\begin{proof}
Every summand is nonnegative for integer arguments. Conversely,
$4x+1$ is not of the form $4^r(8s+7)$, so Legendre's three-square
theorem applies. Exactly one square is odd modulo four. Write the roots
as $2a,2b,2c+1$, subtract one, and divide by four.
\end{proof}

Initially guard all 1,228 remaining semantic natural coordinates using
\eqref{eq:three-square}. The arithmetic wires are natural by topological
induction. This gives
\[
2944+3\cdot1228=6628\quad\text{variables},\qquad
2505+1228=3733\quad\text{rows}.
\]
The next two sections show why most of those guards are unnecessary.

\section{Normalizing signs by exact symmetries}
\label{sec:signs}

A natural witness may be needed only up to a sign that the equations
cannot distinguish. In that case it is enough that every integer
solution has an equivalent representative with the desired sign. This
is a different reason for omitting a guard than proving that the
original coordinate is automatically nonnegative.

\subsection{The general sign criterion}

For a quadratic residual
$f(\boldsymbol X)=\sum_\alpha c_\alpha\boldsymbol X^\alpha$, replace
$X_i$ by $(-1)^{s_i}X_i$, where $s\in\Ftwo^n$. The residual changes
only by an overall sign if
\begin{equation}
\label{eq:parity}
\langle \alpha-\beta,s\rangle=0\pmod2
\quad\text{for every pair of monomials in }\supp(f).
\end{equation}
The coefficients need not be positive. Condition \eqref{eq:parity} is
an exact finite linear system over $\Ftwo$, and it guarantees
$f((-1)^sX)=\pm f(X)$ for every integer assignment. Its intersection
over all residuals is a diagonal sign-symmetry group of the entire
zero set.

\begin{lemma}[Independent sign normalization]
\label{lem:normalize}
Suppose the core system has sign symmetries $\sigma_1,\ldots,\sigma_k$
such that $\sigma_i$ flips a selected semantic coordinate $z_i$ and
changes no other semantic coordinate. Suppose none of the $z_i$ is
guarded. Then every solution of the core plus the remaining semantic
guards can be changed to one with $z_1,\ldots,z_k\ge0$, while preserving
all those guards and all core zero equations.
\end{lemma}
\begin{proof}
Apply $\sigma_i$ exactly when $z_i<0$. It changes $z_i$ to its absolute
value and leaves every other semantic coordinate unchanged. Diagonal
sign changes commute, so later choices do not undo earlier
normalizations. Every core residual is multiplied by a unit $\pm1$,
and each guarded semantic coordinate keeps its value. A zero coordinate
needs no change.
\end{proof}

\subsection{The 110 normalizations in the actual equation}

The core parity matrix has rank 2,834 on 2,944 coordinates. The emitted
certificate supplies 110 independent sign symmetries, two for each of
the 55 power graphs. The checker does not trust a rank calculation as a
semantic premise: it directly checks \eqref{eq:parity} for every
generator and every actual residual.

The first symmetry negates the coordinate $Z$ in the power certificate
and its product wire $WZ$. The product definition is multiplied by
$-1$, while occurrences of its square are unchanged. The second
negates the Pell coordinate $U$ and the two signed congruence quotients
attached to modulus $U$. Products $Uq$ are unchanged, and the Pell
equation uses $U^2$. All other semantic coordinates are fixed.

The first case illustrates why inspecting only the exponents of $Z$
in the final core is insufficient: $Z$ occurs linearly in the
definition of $WZ$. Coordinated sign changes of a semantic witness and
its wires reveal the actual symmetry. The second case is enabled by
the earlier signed-quotient reduction; natural quotient pairs would
obscure the simple sign action.

These 110 normalizations save 330 guard auxiliaries. Some arithmetic
wires may be negative before normalization, so naturality of all wires
must be proved \emph{after} normalization, as in the next section.

\section{Recovering signs from the equations themselves}
\label{sec:bounds}

The remaining guard optimization is an exact proof of lower bounds,
not an assumption that an arithmetic-looking variable is positive.
The proof certificate starts with a chosen seed set of guarded
semantic coordinates and the normalized coordinates from
Lemma~\ref{lem:normalize}. It then derives lower bounds for every
natural core coordinate.

\subsection{A finite positivity calculus}

Maintain a proved lower bound $L_i\in\{0,1,2\}$ for some coordinates;
unbounded coordinates may still have either sign. For a polynomial
$G$, substitute $X_i=L_i+T_i$ for the bounded coordinates, with
$T_i\ge0$. If the resulting coefficients are nonnegative and every
unbounded variable has even exponent in each monomial, then $G\ge0$.
If the constant coefficient is positive, then $G>0$.

This criterion is deliberately sufficient rather than complete. It
can be checked by finite integer polynomial arithmetic. Squares of
unbounded coordinates are allowed, which is useful before every wire
has acquired a sign bound.

\begin{lemma}[Positive-multiplier inference]
\label{lem:bounds}
Suppose an actual residual, oriented by a possible factor $-1$, has
the form
\[
A(\boldsymbol X)x=B(\boldsymbol X),
\]
where $A,B$ do not contain $x$, and previously proved bounds establish
$A>0$. Then:
\begin{enumerate}
\item $B\ge0$ implies $x\ge0$;
\item $B>0$ implies $x\ge1$ for integer $x$;
\item $B-A>0$ implies $x\ge2$ for integer $x$.
\end{enumerate}
\end{lemma}
\begin{proof}
The first two conclusions follow from multiplication by a positive
number. For the third, $A(x-1)=B-A>0$, so $x>1$; integrality gives
$x\ge2$. No division is inserted in the polynomial system.
\end{proof}

Examples include $A=2+s$, which makes the guard on $A$ redundant when
$s\ge0$, and the less obvious row
\[
B=1+4Yq_4.
\]
Once $B\ge2$ and $Y\ge1$ are known, it forces $q_4\ge1$ without a
guard on $q_4$. This is a genuine use of a positive multiplier, not
mere propagation of nonnegative coefficients.

\subsection{The certified seed set}

The complete certificate keeps 645 semantic guards. Its 110 sign
normalizations and 2,382 ordered lower-bound inferences recover
nonnegativity of all 2,723 natural core coordinates. The 221 signed
quotients retain no sign bound. Among the 1,228 semantic natural
coordinates, the exact accounting is
\[
1228=645\ \text{guarded}
\ +\ 110\ \text{normalized}
\ +\ 473\ \text{inferred}.
\]
Some coordinates receive stronger bounds after an initial bound, which
explains why the number of inference steps exceeds the number of new
coordinates proved nonnegative.

The seed selection is an optimization heuristic. Its success is not a
mathematical hypothesis: a separate checker verifies the supplied
ordered proof, including the strict positivity of every multiplier and
the exact polynomial row used at every step. It does not call the
optimizer. Minimality of the 645-coordinate seed set is not asserted.

\begin{theorem}[The 4,879-variable equation]
\label{thm:guarded}
The system consisting of the 2,505-row signed core and the 645 displayed
three-square guards has an integer solution if and only if the frozen
natural parent has a natural solution. Its squared sum has 4,879
integer coordinates, 3,150 quadratic residuals, 22,910 nonzero terms,
and coefficient height 127,285,349.
\end{theorem}
\begin{proof}
For soundness first normalize the 110 signs. The retained guards remain
valid. Replay the lower-bound proof to recover all natural core
coordinates. Apply the signed-quotient lift and the two affine bridge
restorations to obtain a natural parent solution. For completeness,
project any natural parent solution to the signed core. Every one of
its 645 guarded coordinates is natural, so the three-square theorem
provides the fresh guard auxiliaries. All other omitted guards are
unneeded in this direction.

The variable and row counts are
$2944+3\cdot645=4879$ and $2505+645=3150$. Exact expansion of the
supplied finite rows gives the remaining measurements.
\end{proof}

\begin{corollary}[The retained-wire companion]
\label{cor:retained}
The same methods applied to the retained-wire natural parent give an
equivalent integer equation in 7,058 variables, with 5,297 quadratic
rows, 24,277 monomials, and coefficient height 1,992.
\end{corollary}
\begin{proof}
The exact signed-quotient certificate first removes 442 private affine
bridges and the same 221 quotient pairs. The resulting core has
$6180-442-3\cdot221=5075$ variables and
$5520-442-2\cdot221=4636$ rows. It retains 1,228 semantic natural
inputs, 221 signed quotients, and 3,626 arithmetic wires.

The sign certificate uses the same 110 independent semantic
normalizations, now with the additional literal wire sign changes.
It retains 661 guards, 16 more than the affine-source case, and
recovers 457 other semantic signs. Its 4,942 ordered bound steps prove
nonnegativity of all 4,854 natural core coordinates. The same
normalization, inference, and quotient-restoration arguments apply.
Thus the integer dimensions are $5075+3\cdot661=7058$ and
$4636+661=5297$. Full expansion gives the remaining exact measurements.
\end{proof}

\section{Polynomial elimination without a degree increase}
\label{sec:elimination}

After sign normalization and guard reduction, every actual coordinate
of the displayed system ranges over $\Z$. This licenses an especially
simple second reduction: solve monic equations and substitute them.
There is no need to preserve a separate natural domain at this stage,
because its required effects are already encoded by the remaining rows.

\begin{theorem}[Degree-preserving monic projection]
\label{thm:elimination}
Consider a system of polynomials of degree at most two over $\Z$ with
one distinguished equation
\[
z-R(\boldsymbol x)=0,
\]
where $R$ is independent of $z$. Assume either:
\begin{enumerate}
\item $R$ has degree at most one; or
\item $R$ has degree at most two, and every other residual contains
    $z$ only as a scalar-linear term $cz$.
\end{enumerate}
Delete the distinguished row and substitute $R$ for $z$ everywhere
else. The resulting residuals have degree at most two, and projection
off $z$ is a bijection between the two integer zero sets. Its inverse
is the polynomial map $\boldsymbol x\mapsto(\boldsymbol x,R(\boldsymbol x))$.
\end{theorem}
\begin{proof}
Every source solution has $z=R(\boldsymbol x)$ and therefore projects
to a target solution. Conversely, adjoining that value makes the
distinguished row and all substituted rows zero, uniquely. An affine
substitution in a quadratic polynomial remains quadratic. In the
second case each other row is $cz+G(\boldsymbol x)$ and becomes
$cR+G$, also quadratic. All coefficients remain integers, so there is
no divisibility side condition or exceptional denominator.
\end{proof}

\begin{remark}
The second criterion excludes occurrences such as $z^2$ and $zw$.
Substituting a quadratic expression there could raise the residual
degree to four and the squared-sum degree to eight. The implementation
checks actual monomials in every affected row, not just a textual
occurrence count. Solving $2z=R$ without a divisibility certificate is
also excluded; the pivot coefficient must be $1$ or $-1$.
\end{remark}

The algorithm repeatedly chooses a permitted sparse substitution. It
records the exact current defining row, pivot variable, right-hand
side, and full set of affected rows. Original variable and row indices
remain fixed throughout the certificate; surviving coordinates are
renumbered only at the end. A complete reverse lift is obtained by
processing the certificate backwards. Composition of the polynomial
lifts gives an integer polynomial inverse to the final projection on
the zero sets.

The recommended final pass starts with the 7,058-variable equation
of Corollary~\ref{cor:retained} and removes $\ElimVars$ variables:
$\AffineElims$ by affine substitution and $\QuadraticElims$ by the
linear-occurrence quadratic rule. It permits at most two monomials
in each substituted right-hand side, as a size heuristic to avoid
dense coefficient growth. The support caps are not
correctness assumptions. A separate checker imports no optimizer and
replays every exact sparse substitution, including changes to guards.

For the smaller-variable companion, a second pass starts with the
4,879-variable affine-source equation and removes 1,208 coordinates:
461 affine and 747 quadratic substitutions. This gives 3,671 variables,
1,942 residuals, and the larger coefficient height in the comparison
table. That pass uses support caps 18 and 24, respectively. Both complete
equations and both exact inverse-map certificates are retained.

Applying Theorem~\ref{thm:elimination} successively to
Corollary~\ref{cor:retained}, then using
\eqref{eq:parent-correctness}, proves Theorem~\ref{thm:main}. The
same argument applied to Theorem~\ref{thm:guarded} proves the
consistency equivalence for the 3,671-variable companion. The
coefficient differences recorded in Table~\ref{tab:main} are real; a final
polynomial is not identical to its parent polynomial. What is preserved
at this stage is a polynomial bijection of integer zero sets.

\section{Why the equivalence is provable in ZFC}
\label{sec:internal}

There are three distinct assertions that must not be conflated:
\begin{enumerate}
\item a standard-model equivalence checked separately for selected
    finite inputs;
\item an ordinary uniform mathematical proof available within ZFC;
\item a new emitted object-language proof or a complete kernel-checked
    implementation of that proof.
\end{enumerate}
The delivered result is of the second kind, supported by exact
coefficient and transformation certificates. It does not claim the
third kind.

The new transformations are elementary and uniform. A quotient lift
uses integer order and maximum; a sign normalization uses finitely
many polynomial symmetries; the lower-bound proof is a finite sequence
of ordered-ring implications; and each monic elimination is an exact
polynomial identity with a polynomial inverse. ZFC proves the
required laws of its internally defined integers and proves the
three-square theorem. Therefore the finite certificates can be
expanded into ordinary ZFC derivations of these transformations.

The source arithmetic also uses uniform theorems: finite products,
induction, prime factorization, CRT, Pell identities, and syntactic
proof transformations. The carrier and CRT arguments apply to every
internal natural number in a model of ZFC, including numbers that
are nonstandard from an external viewpoint. They are not an appeal to
the assertion that a particular finite diagnostic test happened to pass.
The ordinary C222 calibration and ZF/ZFC refutation transport are
named logical theorems with set-theoretic proofs; their use is
explicit in the dependency ledger.

The consistency sentence is syntactic. The theorem states
\[
\ZFC\vdash
\bigl(\exists\boldsymbol x\in\Z^{\MainVars}\;P(\boldsymbol x)=0\bigr)
\ \longleftrightarrow\
\bigl(\exists p\;\operatorname{Proof}_{\ZFC}(p,\bot)\bigr).
\]
Negating both sides yields the displayed no-zero form. No soundness or
consistency assumption is used to prove the biconditional itself.
In particular, this work does not prove that $P$ lacks solutions.
If ZFC is consistent, a ZFC proof that $P$ lacks solutions would,
through this biconditional, give a ZFC proof of its own consistency.
The construction is compatible with the second incompleteness theorem.

The archive includes ordinary dependency proof notes and pinned source
interfaces, but not a newly rebuilt Lean checkout or a line-by-line
formal ZFC proof. That remaining formalization task should be described
accurately. It is different from leaving an unidentified coefficient,
an arbitrary compiler choice, or a conjectural simulation premise in
the polynomial.

\section{Exact verification and reproducibility}
\label{sec:verification}

The verification is designed around the risks introduced by the new
construction: incorrect occurrence restrictions, sign changes that
affect another guarded variable, circular positivity arguments,
unsafe polynomial substitutions, and a mismatch between the stated
residuals and their expanded squared sum.

The release verifier checks the following distinct objects.
\begin{enumerate}
\item \textbf{The fixed source.} The pinned direct arithmetic generator
    reproduces the complete natural parent. Original independent
    auditors check the actual twelve calculus branches, positive-power
    interfaces, and affine restoration. The package contains the
    relevant mathematical proof notes and licenses.
\item \textbf{Signed quotient blocks.} Every quotient and product use
    is located in the complete source row list. The exact common
    modulus, opposite signs, private-use restrictions, and nonnegative
    restoration of all remaining arithmetic wires are checked.
\item \textbf{Sign symmetries.} Every monomial of every core residual
    satisfies the parity test for each supplied generator. The checker
    also verifies that only the intended semantic coordinate changes.
\item \textbf{Lower-bound derivations.} Each step cites an actual row,
    uses only earlier bounds, and establishes strict positivity of its
    multiplier by exact shifted polynomial arithmetic. At the end,
    every intended natural coordinate has a nonnegative bound.
\item \textbf{Elimination.} A separate implementation reconstructs the
    solved equation and substitutes it into every actual occurrence.
    It verifies the degree criterion, deleted rows, surviving indices,
    and exact final renumbering.
\item \textbf{The polynomial.} All squared residuals are fully expanded
    and collected. Every coefficient, monomial, degree, variable, term
    count, and height is compared to the released file. Hashes identify
    the exact data subjected to these checks.
\end{enumerate}

These checks use exact integer arithmetic. They do not rely on
floating-point agreement, randomly selected zeros, or a finite search
for a contradiction. The checkers use explicit failure conditions;
Python's \texttt{-O} option does not disable their acceptance logic.
Normal and optimized runs are included in the release evidence. The
independent checkers do not import the new optimization routines.

The universal proofs establish why a syntactically valid certificate
preserves existence for all witnesses. The finite checker establishes
that the delivered coefficients and certificates instantiate those
theorems. Neither role replaces the other.

\subsection{Practical commands}

Python 3.11 or later and its standard library suffice for the exact
arithmetic verification. No network call or package installation is
performed by the verification driver.
\begin{lstlisting}
python3 verify.py
python3 -O verify.py
\end{lstlisting}
The README gives the additional regeneration commands, including the
slower optimization pass for both supplied final endpoints. A
standard TeX Live or MiKTeX installation builds this article:
\begin{lstlisting}
cd report
pdflatex -interaction=nonstopmode -halt-on-error report.tex
pdflatex -interaction=nonstopmode -halt-on-error report.tex
\end{lstlisting}
The complete frozen data allow verification independently of later
changes to ProveIt. The intended integration is a new candidate or
incoming research supplement; no existing repository file needs to be
overwritten.

\section{What is new, and what has not been established}

The specific new contribution is the combination of signed-domain
compilation, exact sign symmetries, certified order propagation, and
degree-preserving polynomial elimination on this fixed direct
consistency equation. The first two methods interact: replacing
natural quotient pairs by signed quotients exposes the Pell-modulus
sign symmetry and makes its normalizing map particularly simple.
The sign proof then removes guard variables; unrestricted polynomial
elimination can finally remove semantic slack coordinates as well as
arithmetic wires without maintaining an artificial natural domain.

The constituent mathematical ideas have recognizable precedents:
differences of naturals, sign symmetries, three-square representations,
and monic substitutions are classical. We do not claim that these
general ideas were previously unknown. The exact combined application,
finite certificates, and resulting full polynomial are the advance
relative to the inspected repository snapshot. Establishing priority
in the broader literature would require a separate review.

The following statements are deliberately not made:
\begin{itemize}
\item that this is the fewest possible number of unknowns;
\item that degree four is minimal for this particular consistency
    encoding;
\item that the guarded seed set or elimination order is optimal;
\item that a literal polynomial with fewer variables must have a
    simpler proof or a shorter description;
\item that the finite audits decide ZFC consistency; or
\item that a new complete Lean proof or an explicit formal ZFC proof
    file is present.
\end{itemize}
The useful claim is narrower and fully measured: a large reduction of
the explicit direct quartic's dimension, with the same logical target,
the same exact degree, reproducible ordinary coefficients, and
constructive proof maps.

\section{Further research questions}
\label{sec:future}

\paragraph{1. Optimize guard seeds with a proved lower bound.}
The current seed set is found by a heuristic and certified afterward.
The inference rules define a directed dependency structure with
alternative prerequisites. Can an integer-programming or satisfiability
formulation find a smaller seed set, and can a matching lower-bound
certificate prove optimality within a clearly stated inference calculus?
An improvement here saves three variables per deleted guard before
subsequent eliminations.

\paragraph{2. Use larger sign groups or affine symmetries.}
The current certificate uses independent diagonal sign normalizations.
More general transformations, such as a sign change accompanied by an
integral translation of congruence quotients, may preserve the zero set
even when diagonal parity alone does not. Such a transformation needs
an explicit inverse and must preserve the necessary semantic signs.
Classifying these symmetries of the positive power graph is a concrete
next target.

\paragraph{3. Compile the arithmetic graph directly with mixed domains.}
Here the signed quotient reduction is applied to a frozen natural
parent so that the comparison is exact. A compiler that tracks natural,
strictly positive, sign-normalizable, and unrestricted coordinates from
the outset could avoid creating redundant witnesses. Its output should
still carry a simple independently checked certificate, rather than
require trusting a more elaborate optimizer.

\paragraph{4. Jointly optimize dimension, sparsity, and coefficient height.}
The final monic elimination increases coefficient height. Different
pivot choices or retaining selected scalar wires could yield a better
dimension/height tradeoff. A useful search should maintain the complete
expanded polynomial at each accepted endpoint and report a small
nondominated set rather than promote one arbitrary weighted score.

\paragraph{5. Replace a complete arithmetic primitive.}
The direct construction contains 55 power graphs. A smaller
fully proved positive power representation, or a representation shared
across compatible calls, could yield a larger improvement than further
local row simplification. The exact source/Pell index hypotheses must
be retained. Replacing the graph by a plausible rational Pell condition
without proving integrality is insufficient.

\paragraph{6. Produce the end-to-end internal proof artifact.}
The ordinary ZFC proof is assembled here, but its actual emitted
internal derivation is still a separate deliverable. The new algebraic
and sign certificates are finite and suitable for a small trusted
checker. Formalizing these generic certificate validators once would
make later size optimizations much less expensive to certify. The
source calibration and the new literal coefficient binding should
remain separate theorem interfaces.

\paragraph{7. Compare with small-variable encodings on description length.}
A quartic with thousands of moderate coefficients and a high-degree
ten-variable polynomial optimize different resources. A fair comparison
could measure a fixed binary description format, including all literal
coefficients and any bounded coefficient recipes, together with the
cost of a proof checker. This would make the request for a ``simplest''
equation more precise without confusing arithmetic operation counts
with coefficient bit complexity.

\appendix
\section{The arithmetic criterion behind the sign symmetries}
\label{app:power}

For positive inputs $n,k$ and natural output $m$, the parent power
graph has 27 semantic auxiliary witnesses. Four difference-of-natural
quotients become four signed witnesses after our first reduction,
leaving 23 auxiliary coordinates before wire flattening. Write
\[
\mathcal N(a,x,y)=x^2-(a^2-1)y^2-1.
\]
The full positive criterion consists of the following seventeen
constraints. All auxiliary coordinates are initially natural except
the four new congruence quotients $z_1,z_2,z_3,z_4$:
\begin{align*}
A&=a_2+2,&Y&=k+y_k,&B&=b_2+2,\\
\mathcal N(A,X,Y)&=0,&\mathcal N(A,U,V)&=0,&\mathcal N(B,S,R)&=0,\\
B&=1+4Yq_4,&V&=Y^2(v_d+1),\\
B+Uz_1&=A,&S+Uz_2&=X,\\
R+4Yz_3&=k,&X+Yn+Tz_4&=YA+m,\\
2An&=T+n^2+1,&T&=m+g+1,\\
W&=n+w_n,&W&=k+w_k,\\
\mathcal N(W+1,A,WZ)&=0.&&
\end{align*}
These expressions are displayed before arithmetic wire flattening;
their degrees are not the degree bound on the final residuals. The
actual generator replaces their product expressions by the retained
quadratic gate equations and shares repeated subexpressions.

The last norm rewrites as
\[
A^2=1+W^3(W+2)Z^2.
\]
Since $A\ge2$ and $W\ge n\ge1$, it forces $Z\ne0$ and
$A>W\ge n$. Therefore the natural truncated difference
$A\mathbin{\dot-}n$ in the inherited power theorem equals the ordinary
integer difference $A-n$ in the rearranged congruence. Also
$Y\ge k\ge1$, so $V=Y^2(v_d+1)>0$ and the zero-index Pell branch is
absent. These facts justify the positive specialization of the
imported power criterion.

The sign symmetries are now visible at the unflattened level:
$Z\mapsto-Z$ changes no norm, and
$(U,z_1,z_2)\mapsto(-U,-z_1,-z_2)$ changes neither the norms nor the
congruences. The literal flattened certificate checks every wire
needed to realize the same action. Bounds inferred from these rows
are used only after the necessary natural coordinates are guarded or
normalized, so the proof does not silently assume $V>0$ when
$v_d=-1$ might otherwise be possible.

\section{Data and proof-certificate conventions}

Variable lists are ordered and indexed from zero. Each sparse polynomial
uses sorted monomials and nonzero integer coefficients; duplicates are
collected. The final expanded table has exactly $\MainTerms$ entries.
The main equation and its expansion are identified by the hashes in
\texttt{equations/main/BUILD.json} and the release manifest. The
polynomial hash is
\begin{center}\small\ttfamily\MainPolynomialHash\end{center}

For an elimination step, the certificate stores an original coordinate
index $v$, a current monic defining row, a polynomial expression $R$,
and the full set of current rows containing $v$. The checker first
verifies the identity of the defining row with $\pm(X_v-R)$, then
checks the degree criterion and substitutes into every occurrence.
The certificate can be interpreted backwards without rerunning the
optimization search.

For a sign step, the certificate stores the list of flipped coordinate
indices and the unique semantic pivot. For a lower-bound step it stores
the target coordinate, new bound, and the actual source row index.
The multiplier and right-hand side are reconstructed from that row by
the checker. Therefore a claimed lower bound cannot smuggle an
additional equation into the proof.

\section{A small illustrative quotient block}

This example is explanatory and does not replace the full C222 instance.
Suppose a retained natural modulus is $M=7$ and a row requires a
difference $E=-21$. The old equations
\[
u=7q_+,\quad v=7q_-,\quad -21+u-v=0
\]
can be replaced by $-21+7z=0$. The signed witness $z=3$ restores as
$q_+=3,q_-=0,u=21,v=0$. Reversing the sign of $E$ gives $z=-3$ and
restores with the two natural quotient parts exchanged. If $M=0$,
both systems require only $E=0$, and the canonical lift still works.
This last case explains why avoiding a division by $M$ is valuable.

For an illustrative elimination, the rows
\[
z-x^2-y^2=0,\qquad 3z+w-7=0
\]
reduce to $3x^2+3y^2+w-7=0$, with inverse $z=x^2+y^2$ and no degree
increase. If the second row instead contains $z^2$, this quadratic
elimination is disallowed by the certificate rule.

\section{Provenance and scope of the release}

The comparison snapshot is fixed even if the repository advances while
this research is in progress. The archive retains the exact natural
source, original arithmetic generator and auditors, relevant ordinary
proof notes, source-calibration interfaces, licenses, the new complete
integer equations, independent checkers, and the article source and
compiled PDF. Intermediate experimental outputs that do not support a
released endpoint are omitted.

No remote repository mutation is needed to use the deliverable. The
package can be reviewed and placed under
\texttt{Computability/DiophantineConsistency/candidates} or
\texttt{docs/incoming} at the owner's discretion. The release manifest
records the exact file bytes; reproducing a PDF on a different TeX
installation may legitimately change PDF metadata while leaving its
mathematical content unchanged.

\begin{thebibliography}{20}
\bibitem{repo-integer}
V. Reshetnikov and contributors, \emph{An explicit quartic over 8,619
unrestricted integers}, ProveIt, pinned snapshot
\texttt{278bc68df1af125d62f9cbfe28c9fd3ebcf50d5b},
\url{https://github.com/VladimirReshetnikov/ProveIt/blob/278bc68df1af125d62f9cbfe28c9fd3ebcf50d5b/Computability/DiophantineConsistency/candidates/shared-crt-prime-lift/integer/README.md}.

\bibitem{repo-prime}
V. Reshetnikov and contributors, \emph{Factorial-free CRT without a
factorial-size bound}, same snapshot,
\url{https://github.com/VladimirReshetnikov/ProveIt/blob/278bc68df1af125d62f9cbfe28c9fd3ebcf50d5b/Computability/DiophantineConsistency/candidates/shared-crt-prime-lift/PROOF.md}.

\bibitem{repo-direct-review}
V. Reshetnikov and contributors, \emph{Ordinary correctness review:
direct C222 quartic, round 12}, same snapshot,
\url{https://github.com/VladimirReshetnikov/ProveIt/blob/278bc68df1af125d62f9cbfe28c9fd3ebcf50d5b/Computability/DiophantineConsistency/PROPOSAL_REVIEW_DIRECT_C222_ROUND12.md}.

\bibitem{repo-shared-review}
V. Reshetnikov and contributors, \emph{Shared-direct certificate
review}, same snapshot,
\url{https://github.com/VladimirReshetnikov/ProveIt/blob/278bc68df1af125d62f9cbfe28c9fd3ebcf50d5b/Computability/DiophantineConsistency/PROPOSAL_REVIEW_SHARED_CERTIFICATES_DIRECT.md}.

\bibitem{repo-frontier}
V. Reshetnikov and contributors, \emph{Later arrivals and current size
reductions}, same snapshot,
\url{https://github.com/VladimirReshetnikov/ProveIt/blob/278bc68df1af125d62f9cbfe28c9fd3ebcf50d5b/Computability/DiophantineConsistency/PROPOSAL_REVIEW_LATEST.md}.

\bibitem{sun2026}
Z.-W. Sun, \emph{Ten unknowns for Hilbert's tenth problem over the
integers}, arXiv:2609.01594v2, 2026 preprint,
\url{https://arxiv.org/html/2609.01594v2}.

\bibitem{sun2021}
Z.-W. Sun, \emph{Further results on Hilbert's tenth problem},
Science China Mathematics \textbf{64} (2021), 281--306;
\url{https://arxiv.org/abs/1704.03504}.

\bibitem{carl-moroz}
M. Carl and B. Z. Moroz, \emph{On a Diophantine representation of the
predicate of provability}, MPIM Preprint Series 2007 (132),
\url{https://archive.mpim-bonn.mpg.de/id/eprint/581/}.

\bibitem{bayer-david}
J. Bayer and M. David, \emph{A Formal Proof of Complexity Bounds on
Diophantine Equations}, ITP 2025, LIPIcs \textbf{352}, 3:1--3:18,
\url{https://doi.org/10.4230/LIPIcs.ITP.2025.3}.
\end{thebibliography}
\end{document}
